%% kapa: exam=1 round=30 year=2019 period=2 form=A subject=appraisal_law subject_ko=감정평가관계법규 seq=34 number=34
%% source: Q-Net 감정평가사 기출문제 (cst00309) artlSeq=5205205 file=extracted/1008/5205205/2019년 제30회 감정평가사 1차시험 문제지/2019년 제30회 감정평가사 1차시험 문제지/2019년도 제30회 감정평가사 1차 2교시 A형.hwp
%% conversion: kordoc 4.24.51 -> tools/gen1.py (auto); review=proofread-sonnet; flags=none
\begin{question}{34}
부동산등기법령상 등기부에 관한 설명으로 옳은 것은?
\begin{choices}
  \item 등기부는 토지등기부, 건물등기부, 집합건물등기부로 구분한다.
  \item 등기부와 폐쇄한 등기기록은 모두 영구히 보존하여야 한다.
  \item 등기부는 법관이 발부한 영장에 의하여 압수하는 경우 외에는 등기정보중앙관리소에 보관ㆍ관리하여야 한다.
  \item 등기기록의 열람 청구는 관할 등기소가 아닌 등기소에 대하여 할 수 없다.
  \item 1동의 건물을 구분한 건물에 있어서는 1동의 건물에 속하는 전부에 대하여 1개의 등기기록을 사용할 수 없다.
\end{choices}
\end{question}
\answer{2}
